\addcontentsline{toc}{section}{СПИСОК ИСПОЛЬЗОВАННЫХ ИСТОЧНИКОВ}

\begin{thebibliography}{99}

	\bibitem{gost19102} ГОСТ 19.102--77. Единая система программной документации. Стадии разработки : межгосударственный стандарт : издание официальное : утвержден Постановлением Государственного комитета стандартов Совета Министров СССР от 20 декабря 1977~г. № 2851 : дата введения 1980-01-01. -- Москва~: Стандартинформ, 2010. -- Текст~: непосредственный.

	\bibitem{gost34601} ГОСТ 34.601--90. Информационная технология. Комплекс стандартов на автоматизированные системы. Автоматизированные системы. Стадии создания : межгосударственный стандарт : издание официальное : утвержден и введен в действие Постановлением Государственного комитета СССР по управлению качеством продукции и стандартам от 29 декабря 1990~г. № 3469 : дата введения 1992-01-01. -- Москва~: Стандартинформ, 2009. -- Текст~: непосредственный.

	\bibitem{gost705} ГОСТ Р 7.0.5--2008. Система стандартов по информации, библиотечному и издательскому делу. Библиографическая ссылка. Общие требования и правила составления : национальный стандарт Российской Федерации : издание официальное : утвержден и введен в действие Приказом Федерального агентства по техническому регулированию и метрологии от 28 апреля 2008~г. № 95-ст : дата введения 2009-01-01. -- Москва~: Стандартинформ, 2008. -- Текст~: непосредственный.

	\bibitem{gost25010} ГОСТ Р ИСО/МЭК 25010--2015. Информационные технологии. Системная и программная инженерия. Требования и оценка качества систем и программного обеспечения (SQuaRE). Модели качества систем и программных продуктов : национальный стандарт Российской Федерации : издание официальное : утвержден и введен в действие Приказом Федерального агентства по техническому регулированию и метрологии от 29 мая 2015~г. № 442-ст : дата введения 2016-06-01. -- Москва~: Стандартинформ, 2015. -- Текст~: непосредственный.

	\bibitem{gost21500} ГОСТ Р ИСО 21500--2023. Управление проектами, программами и портфелями проектов. Контекст и основные понятия : национальный стандарт Российской Федерации : издание официальное : утвержден и введен в действие Приказом Федерального агентства по техническому регулированию и метрологии от 30 октября 2023~г. № 1296-ст : дата введения 2024-06-01. -- Москва~: Российский институт стандартизации, 2023. -- Текст~: непосредственный.

	\bibitem{gost54869} ГОСТ Р 54869--2011. Проектный менеджмент. Требования к управлению проектом : национальный стандарт Российской Федерации : издание официальное : утвержден и введен в действие Приказом Федерального агентства по техническому регулированию и метрологии от 22 декабря 2011~г. № 1582-ст : дата введения 2012-09-01. -- Москва~: Стандартинформ, 2012. -- Текст~: непосредственный.

	\bibitem{omg_uml} OMG Unified Modeling Language (OMG UML). Version 2.5.1 : specification / Object Management Group. -- Needham~: OMG, 2017. -- 796 p. -- URL: https://www.omg.org/spec/UML/2.5.1 (дата обращения: 10.05.2026).
	\bibitem{anderson_kanban} Андерсон, Д. Канбан. Альтернативный путь в Agile / Д. Андерсон~; перевод с английского А. Коробейникова. -- Москва~: Манн, Иванов и Фербер, 2017. -- 336 с. -- ISBN 978-5-00100-530-8. -- Текст~: непосредственный.

	\bibitem{akhmadullin_pms} Ахмадуллин, Д. Ф. Системы управления проектами: анализ существующих программных решений / Д. Ф. Ахмадуллин. -- DOI 10.24411/2658-4964-2020-1120. // StudNet. -- 2020. -- №~9. -- URL: https://cyberleninka.ru/article/n/sistemy-upravleniya-proektami-analiz-suschestvuyuschih-programmnyh-resheniy (дата обращения: 10.05.2026).

	\bibitem{art_gantt} Богнюков, А. А. Применение инструментов проектного планирования: диаграмма Ганта и сетевой график / А. А. Богнюков, Д. Ю. Зорькин, И. А. Тарасова. // Известия Южного федерального университета. Технические науки. -- 2025. -- №~4. -- С. 102--110. -- URL: https://cyberleninka.ru/article/n/primenenie-instrumentov-proektnogo-planirovaniya-diagramma-ganta-i-setevoy-grafik (дата обращения: 10.05.2026).

	\bibitem{beck_xp} Бек, К. Экстремальное программирование. Разработка через тестирование / К. Бек~; перевод с английского. -- Санкт-Петербург~: Питер, 2017. -- 224 с. -- (Библиотека программиста). -- ISBN 978-5-496-02570-6. -- Текст~: непосредственный.

	\bibitem{booch_uml_guide} Буч, Г. Язык UML. Руководство пользователя / Г. Буч, Дж. Рамбо, И. Якобсон~; перевод с английского. -- 2-е изд. -- Москва~: ДМК Пресс, 2015. -- 496 с. -- ISBN 978-5-94074-644-7. -- Текст~: непосредственный.

	\bibitem{booch_ooa} Буч, Г. Объектно-ориентированный анализ и проектирование с примерами приложений / Г. Буч, Р. А. Максимчук, М. У. Энгл [и др.]~; перевод с английского. -- 3-е изд. -- Москва~: ООО <<И.Д. Вильямс>>, 2010. -- 720 с. -- ISBN 978-5-8459-1401-9. -- Текст~: непосредственный.

	\bibitem{art_python_web} Васильев, П. А. Web-программирование на языке Python. Фреймворки Django, Flask / П. А. Васильев. // Наука, техника и образование. -- 2016. -- №~12 (30). -- URL: https://cyberleninka.ru/article/n/web-programmirovanie-na-yazyke-python-freymvorki-django-flask (дата обращения: 10.05.2026).

	\bibitem{art_client_server} Гайнанова, Р. Ш. Создание клиент-серверных приложений / Р. Ш. Гайнанова, О. А. Широкова. // Вестник Казанского технологического университета. -- 2017. -- Т.~20, №~17. -- URL: https://cyberleninka.ru/article/n/sozdanie-klient-servernyh-prilozheniy (дата обращения: 10.05.2026).

	\bibitem{gamma_patterns} Гамма, Э. Паттерны объектно-ориентированного проектирования / Э. Гамма, Р. Хелм, Р. Джонсон, Дж. Влиссидес~; перевод с английского. -- Санкт-Петербург~: Питер, 2020. -- 448 с. -- (Библиотека программиста). -- ISBN 978-5-4461-1595-2. -- Текст~: непосредственный.

	\bibitem{date_db} Дейт, К. Дж. Введение в системы баз данных / К. Дж. Дейт~; перевод с английского. -- 8-е изд. -- Москва~: ООО <<И.Д. Вильямс>>, 2005. -- 1328 с. -- ISBN 5-8459-0788-8. -- Текст~: непосредственный.

	\bibitem{art_scrum_principles} Дружинина, В. Д. Основные принципы методологии SCRUM для управления проектами / В. Д. Дружинина. // Экономика. Социология. Право. -- 2022. -- №~1 (25). -- URL: https://cyberleninka.ru/article/n/osnovnye-printsipy-metodologii-scrum-dlya-upravleniya-proektami (дата обращения: 10.05.2026).

	\bibitem{art_agile_practice} Зайцева, И. А. Практика применения методологий Agile, Scrum в ИТ-проектах / И. А. Зайцева, В. Ш. Ебата, Н. А. Ковбаса. // Индустриальная экономика. -- 2021. -- №~1. -- URL: https://cyberleninka.ru/article/n/praktika-primeneniya-metodologiy-agile-scrum-v-it-proektah (дата обращения: 10.05.2026).

	\bibitem{art_db_normalization} Игнатюк, В. А. Логическое проектирование реляционной базы данных на основе принципов нормализации / В. А. Игнатюк, Е. А. Сторожок. // Территория новых возможностей. Вестник Владивостокского государственного университета экономики и сервиса. -- 2011. -- №~4 (13). -- URL: https://cyberleninka.ru/article/n/logicheskoe-proektirovanie-relyatsionnoy-bazy-dannyh-na-osnove-printsipov-normalizatsii (дата обращения: 10.05.2026).

	\bibitem{art_hashing} Айдинян, А. Р. Эффективность применения некоторых алгоритмов хеширования для веб-приложения с интерфейсами регистрации и аутентификации пользователей / А. Р. Айдинян, О. Л. Легонько, Н. С. Мамонов. // Инженерный вестник Дона. -- 2025. -- №~8. -- URL: https://cyberleninka.ru/article/n/effektivnost-primeneniya-nekotoryh-algoritmov-heshirovaniya-dlya-veb-prilozheniya-s-interfeysami-registratsii-i-autentifikatsii (дата обращения: 10.05.2026).

	\bibitem{art_rest} Ковалев, В. В. Актуальность использования архитектуры REST для обмена данными между клиент-серверными приложениями / В. В. Ковалев, В. В. Храмков, Е. И. Семенова. // Актуальные проблемы авиации и космонавтики. -- 2019. -- Т.~2. -- URL: https://cyberleninka.ru/article/n/aktualnost-ispolzovaniya-arhitektury-rest-dlya-obmena-dannymi-mezhdu-klient-servernymi-prilozheniyami (дата обращения: 10.05.2026).

	\bibitem{connolly_db} Коннолли, Т. Базы данных. Проектирование, реализация и сопровождение. Теория и практика / Т. Коннолли, К. Бегг~; перевод с английского. -- 3-е изд. -- Москва~: ООО <<И.Д. Вильямс>>, 2017. -- 1440 с. -- ISBN 978-5-8459-2020-1. -- Текст~: непосредственный.

	\bibitem{art_digital_pm} Куренков, А. Л. Управление проектами цифровой трансформации: современные тенденции и особенности / А. Л. Куренков. // Вестник Российского экономического университета им.~Г.~В.~Плеханова. -- 2024. -- №~5. -- URL: https://cyberleninka.ru/article/n/upravlenie-proektami-tsifrovoy-transformatsii-sovremennye-tendentsii-i-osobennosti (дата обращения: 10.05.2026).

	\bibitem{laure_api} Лоре, А. Проектирование веб-API / А. Лоре~; перевод с английского. -- Москва~: ДМК Пресс, 2020. -- 440 с. -- ISBN 978-5-97060-861-6. -- Текст~: непосредственный.

	\bibitem{lutz_python} Лутц, М. Изучаем Python / М. Лутц~; перевод с английского. -- 5-е изд. -- Санкт-Петербург~: ООО <<Диалектика>>, 2019--2020. -- Т.~1: 832 с. -- ISBN 978-5-907144-52-1. -- Текст~: непосредственный.

	\bibitem{mazur2014} Управление проектами : учебное пособие / И. И. Мазур, В. Д. Шапиро, Н. Г. Ольдерогге, А. В. Полковников~; под общей редакцией И. И. Мазура и В. Д. Шапиро. -- 10-е изд., стер. -- Москва~: Омега-Л, 2014. -- 960 с. -- ISBN 978-5-370-02800-7. -- Текст~: непосредственный.

	\bibitem{art_uml_modeling} Макеева, О. В. Моделирование информационных процессов с помощью UML / О. В. Макеева, М. В. Сартаков, Е. А. Чернов. // Инновации и инвестиции. -- 2021. -- №~5. -- URL: https://cyberleninka.ru/article/n/modelirovanie-informatsionnyh-protsessov-s-pomoschyu-uml (дата обращения: 10.05.2026).

	\bibitem{art_agile_modern} Мокшин, В. В. Современные подходы к проектам. Методологии. AGILE: SCRUM / В. В. Мокшин, А. М. Гайнутдинова, С. О. Самсонов. // StudNet. -- 2021. -- №~6. -- URL: https://cyberleninka.ru/article/n/sovremennye-podhody-k-proektam-metodologii-agile-scrum (дата обращения: 10.05.2026).

	\bibitem{orlov_se} Орлов, С. А. Программная инженерия. Технологии разработки программного обеспечения : учебник для вузов / С. А. Орлов. -- 5-е изд., обнов. и доп. -- Санкт-Петербург~: Питер, 2020. -- 640 с. -- (Учебник для вузов. Стандарт третьего поколения). -- ISBN 978-5-4461-1348-4. -- Текст~: непосредственный.

	\bibitem{art_db_methodology} Понин, Ф. Н. Методология проектирования и создания баз данных для современного программного обеспечения / Ф. Н. Понин. // Universum: технические науки. -- 2024. -- №~1 (118). -- URL: https://cyberleninka.ru/article/n/metodologiya-proektirovaniya-i-sozdaniya-baz-dannyh-dlya-sovremennogo-programmnogo-obespecheniya (дата обращения: 10.05.2026).

	\bibitem{art_orm} Романов, С. С. Достоинства, недостатки и альтернативы объектно-реляционного отображения (ORM) / С. С. Романов. // Таврический научный обозреватель. -- 2016. -- №~11-2 (16). -- URL: https://cyberleninka.ru/article/n/dostoinstva-nedostatki-i-alternativy-obektno-relyatsionnogo-otobrazheniya-orm (дата обращения: 10.05.2026).

	\bibitem{robson_html} Робсон, Э. Изучаем HTML и CSS / Э. Робсон, Э. Фримен~; перевод с английского. -- 2-е изд. -- Санкт-Петербург~: Питер, 2014. -- 720 с. -- (Head First O'Reilly). -- ISBN 978-5-496-00653-8. -- Текст~: непосредственный.

	\bibitem{sutherland_scrum} Сазерленд, Дж. Scrum. Революционный метод управления проектами / Дж. Сазерленд~; перевод с английского. -- Москва~: Манн, Иванов и Фербер, 2016. -- 272 с. -- ISBN 978-5-00057-722-6. -- Текст~: непосредственный.

	\bibitem{safonova_is} Сафонова, А. А. Информационные системы управления проектами / А. А. Сафонова, О. Н. Куксачева. // Формула менеджмента. -- 2020. -- №~1. -- URL: https://cyberleninka.ru/article/n/informatsionnye-sistemy-upravleniya-proektami (дата обращения: 10.05.2026).

	\bibitem{art_dbms_compare} Басов, А. С. Сравнение современных СУБД / А. С. Басов. // Вестник науки. -- 2020. -- №~9 (30). -- URL: https://cyberleninka.ru/article/n/sravnenie-sovremennyh-subd (дата обращения: 10.05.2026).

	\bibitem{terehova_review} Терехова, А. Е. Обзор методологий управления проектами / А. Е. Терехова, Н. Ю. Верба. // Вестник университета. -- 2014. -- №~2. -- URL: https://cyberleninka.ru/article/n/obzor-metodologiy-upravleniya-proektami (дата обращения: 10.05.2026).

	\bibitem{art_web_security} Усков, И. Д. Безопасность веб-приложений: XSS, CSRF и современные методы защиты / И. Д. Усков, В. Д. Сандаков. // Вестник науки. -- 2026. -- URL: https://cyberleninka.ru/article/n/bezopasnost-veb-prilozheniy-xss-csrf-i-sovremennye-metody-zaschity (дата обращения: 10.05.2026).

	\bibitem{fowler_uml} Фаулер, М. UML. Основы : краткое руководство по стандартному языку объектного моделирования / М. Фаулер~; перевод с английского. -- 3-е изд. -- Санкт-Петербург~: Символ-Плюс, 2013. -- 192 с. -- ISBN 5-93286-060-X. -- Текст~: непосредственный.

	\bibitem{fowler_patterns} Фаулер, М. Шаблоны корпоративных приложений / М. Фаулер, Д. Райс, М. Фоммел [и др.]~; перевод с английского. -- Москва~: ООО <<И.Д. Вильямс>>, 2018. -- 544 с. -- ISBN 978-5-8459-1611-2. -- Текст~: непосредственный.

	\bibitem{flanagan_js} Флэнаган, Д. JavaScript. Полное руководство / Д. Флэнаган~; перевод с английского. -- 7-е изд. -- Москва~: ООО <<Диалектика>>, 2021. -- 720 с. -- ISBN 978-5-907203-79-2. -- Текст~: непосредственный.

	\bibitem{forcier_django} Форсье, Дж. Django. Разработка веб-приложений на Python / Дж. Форсье, П. Биссекс, У. Чан~; перевод с английского. -- Санкт-Петербург~: Символ-Плюс, 2009. -- 456 с. -- ISBN 978-5-93286-167-7. -- Текст~: непосредственный.

	\bibitem{chinnathambi_react} Чиннатамби, К. Изучаем React / К. Чиннатамби~; перевод с английского. -- 2-е изд. -- Москва~: Эксмо, 2019. -- 368 с. -- (Мировой компьютерный бестселлер). -- ISBN 978-5-04-098028-4. -- Текст~: непосредственный.

	\bibitem{art_api_compare} Шор, А. М. Сравнительный анализ подходов в разработке API веб-приложений / А. М. Шор. // StudNet. -- 2020. -- №~9. -- URL: https://cyberleninka.ru/article/n/sravnitelnyy-analiz-podhodov-v-razrabotke-api-veb-prilozheniy (дата обращения: 10.05.2026).

	\bibitem{art_hybrid_pm} Яковчук, А. И. Гибридный подход к управлению проектами планирования мероприятий / А. И. Яковчук. // Экономический вектор. -- 2022. -- Т.~28, №~1. -- С.~56--60. -- URL: https://cyberleninka.ru/article/n/gibridnyy-podhod-k-upravleniyu-proektami-planirovaniya-meropriyatiy (дата обращения: 10.05.2026).

	\bibitem{jacobson_rup} Якобсон, А. Унифицированный процесс разработки программного обеспечения / А. Якобсон, Г. Буч, Дж. Рамбо~; перевод с английского. -- Санкт-Петербург~: Питер, 2002. -- 496 с. -- (Классика computer science). -- ISBN 5-318-00358-3. -- Текст~: непосредственный.

	\bibitem{pmbok6} Руководство к своду знаний по управлению проектами (Руководство PMBOK). Шестое издание + Agile : практическое руководство / Project Management Institute~; перевод с английского. -- Москва~: Олимп-Бизнес, 2022. -- 974 с. -- ISBN 978-5-9693-0402-4. -- Текст~: непосредственный.

	\bibitem{banks_react} Banks, A. Learning React. Modern Patterns for Developing React Apps / A. Banks, E. Porcello. -- 2nd ed. -- Sebastopol~: O'Reilly Media, 2020. -- 310 p. -- ISBN 978-1-492-05172-5. -- Текст~: непосредственный.

	\bibitem{django_docs} Django documentation : официальный сайт. -- San Francisco~: Django Software Foundation, 2005--. -- URL: https://docs.djangoproject.com (дата обращения: 10.05.2026).
	\bibitem{drf_docs} Django REST framework : сайт. -- 2011--. -- URL: https://www.django-rest-framework.org (дата обращения: 10.05.2026).
	\bibitem{mdn_docs} MDN Web Docs : официальный сайт. -- Mountain View~: Mozilla Foundation, 2005--. -- URL: https://developer.mozilla.org (дата обращения: 10.05.2026).
	\bibitem{owasp_top10} OWASP Top 10:2021 // The OWASP Foundation~: [сайт]. -- 2021. -- URL: https://owasp.org/Top10/2021/ (дата обращения: 10.05.2026).
	\bibitem{pmi_standards} PMI Standards \& Publications // Project Management Institute~: [сайт]. -- Newtown Square, 1969--. -- URL: https://www.pmi.org/standards (дата обращения: 10.05.2026).
	\bibitem{react_docs} React : официальный сайт. -- Menlo Park~: Meta Platforms, 2013--. -- URL: https://react.dev (дата обращения: 10.05.2026).
	\bibitem{rfc9110} Hypertext Transfer Protocol (HTTP) Semantics : RFC 9110 / R. Fielding, M. Nottingham, J. Reschke. -- Internet Engineering Task Force, June 2022. -- URL: https://datatracker.ietf.org/doc/html/rfc9110 (дата обращения: 10.05.2026).
	\bibitem{yougile_review} Лучшие таск-трекеры для управления проектами и задачами в 2026 году: обзор 19 российских систем. // Хабр. -- 2025. -- URL: https://habr.com/ru/companies/yougile/articles/991204/ (дата обращения: 10.05.2026).

\end{thebibliography}
